\section{Chapter Workshop}

\begin{workedexample}{Least squares as orthogonal projection}
Fit a line through the origin, $y=ax$, to the data $(1,1)$, $(2,2)$, and $(3,2)$. In matrix form,
\[
    A=\begin{pmatrix}1\\2\\3\end{pmatrix},
    \qquad
    b=\begin{pmatrix}1\\2\\2\end{pmatrix}.
\]
The least-squares coefficient satisfies the normal equation
\[
    A^{T}A\,a=A^{T}b,
\]
so
\[
    14a=11,
    \qquad a=\frac{11}{14}.
\]
The residual $r=b-Aa$ is orthogonal to the column space of $A$ because $A^Tr=0$. The computation is therefore a projection, not merely a formula for fitting data.
\end{workedexample}

\begin{applicationbox}{Principal component analysis}
After centering a data matrix $X$, the eigenvectors of $X^TX$ give directions of maximal sample variance. Equivalently, the right singular vectors of $X$ provide orthogonal coordinates ordered by singular value. Truncating the singular-value decomposition produces a low-rank approximation and is used in compression, denoising, visualization, and latent-feature extraction.
\end{applicationbox}

\begin{chapterexercises}
    \item Row-reduce $\begin{pmatrix}1&2&1\\2&4&0\\1&2&3\end{pmatrix}$ and find its rank and nullity.
    \item Find an orthonormal basis for the span of $(1,1,0)$ and $(1,0,1)$.
    \item Determine the eigenvalues and eigenvectors of $\begin{pmatrix}2&1\\1&2\end{pmatrix}$.
    \item Prove that eigenvectors corresponding to distinct eigenvalues of a real symmetric matrix are orthogonal.
    \item Find the least-squares solution of $x\approx1$, $x\approx2$, $x\approx6$.
    \item Explain why a matrix and the linear map it represents are not literally the same object.
    \item Find a basis for the column space and null space of $\begin{pmatrix}1&0&2\\0&1&3\\0&1&3\end{pmatrix}$.
    \item Compute the standard matrix of the linear transformation $T(x,y)=(x+y,\,x-y)$.
    \item Decide whether $\{(1,0,0),(0,1,0),(1,1,0)\}$ is linearly independent.
    \item Diagonalize $\begin{pmatrix}3&0\\0&1\end{pmatrix}$ and interpret its eigenvalues.
    \item Use Gram-Schmidt to orthogonalize $(1,1)$ and $(1,0)$ in $\mathbb{R}^2$.
    \item Verify directly that $\operatorname{rank}(A)+\operatorname{nullity}(A)=n$ for $A=\begin{pmatrix}1&2\\2&4\end{pmatrix}$.
\end{chapterexercises}

\begin{selectedhints}
    \item The second column is twice the first; elimination gives rank $2$ and nullity $1$.
    \item Normalize the first vector, subtract its projection from the second, then normalize.
    \item The eigenvalues are $3$ and $1$, with directions $(1,1)$ and $(1,-1)$.
    \item Compare $\langle Av,w\rangle$ with $\langle v,Aw\rangle$.
    \item Minimize $(x-1)^2+(x-2)^2+(x-6)^2$; the answer is the mean, $x=3$.
    \item A matrix depends on chosen bases; the abstract map does not. A change of basis changes the matrix by similarity or by left/right coordinate transformations.
    \item The pivot columns form a basis for $\operatorname{Col}(A)$; solve $A\mathbf{x}=\mathbf{0}$ for $\operatorname{Nul}(A)$.
    \item Apply $T$ to the standard basis vectors $e_1=(1,0)$ and $e_2=(0,1)$.
    \item Check whether one vector is a linear combination of the others, or row-reduce the matrix with these vectors as columns.
    \item The matrix is already diagonal, so the eigenvalues are its diagonal entries and the eigenvectors are the coordinate axes.
    \item Subtract the projection of the second vector onto the first, then normalize the resulting orthogonal pair.
    \item This matrix has one pivot column, so its rank is $1$ and its nullity is $1$.
\end{selectedhints}
